\documentclass[reqno,11pt]{amsart}
\usepackage{math327-lecture}
\usepackage{needspace}
\usepackage{tikz}
\lectureinfo{18}{Oct. 21, 2026}{Dihedral Groups, the Correspondence Theorem, and Direct Products}



\tcolorboxenvironment{example}{lecturebox,colback=definitiongreen,colframe=definitionborder}
\definecolor{remarkpurple}{HTML}{F2ECFA}
\definecolor{remarkpurpleborder}{HTML}{A68CC7}
\tcolorboxenvironment{remark}{lecturebox,colback=remarkpurple,colframe=remarkpurpleborder}

\begin{document}
\raggedbottom
\makelecturetitle

\section{Dihedral groups}

\begin{define}[Dihedral group]
	For an integer $n\geq3$, the \emph{dihedral group} $D_{2n}$ is the group generated by a rotation $x$ of
	order $n$ and a reflection $y$ of order $2$, with relations
	\[
	x^n=1,\qquad y^2=1,\qquad yxy^{-1}=x^{-1}.
	\]
The last relation is equivalent to $yx=x^{-1}y=x^{n-1}y$.
Our notation records the order of the group in its subscript. One can verify that the relation $yx=x^{-1}y$ allows one to write every element as
either $x^i$ or $x^iy$, where $0\leq i<n$. So $D_{2n}$ consists of 
$2n$ elements:
\[
\{1,x,x^2,\ldots,x^{n-1},y,xy,x^2y,\ldots,x^{n-1}y\}.
\]
\end{define}

\begin{example}
	
	Last time we dealt with $n=4$ and the square. Consider now the regular hexagon, $n=6$ and $D_{12}$. It has 12 elements. What are they ? Label the vertices counterclockwise as shown below.
The blue arrows show two rotations; the dashed lines show two
reflection axes.
\begin{center}
\begin{tikzpicture}[x=1cm,y=1cm,every node/.style={font=\small,align=center}]
  \foreach \offset/\panel in {0/1,3.5/2,7/3,10.5/4} {
    \begin{scope}[xshift=\offset cm]
      \ifnum\panel=1
        \draw[-stealth,courseblue,line width=1pt]
          (0:1.6) arc[start angle=0,end angle=60,radius=1.6];
      \fi
      \ifnum\panel=2
        \draw[-stealth,courseblue,line width=1pt]
          (0:1.6) arc[start angle=0,end angle=120,radius=1.6];
      \fi
      \ifnum\panel=3
        \draw[dashed,ForestGreen,line width=1pt] (180:1.6)--(0:1.6);
      \fi
      \ifnum\panel=4
        \draw[dashed,BrickRed,line width=1pt] (210:1.6)--(30:1.6);
      \fi
      \foreach \vertex in {1,...,6} {
        \coordinate (v\vertex) at ({60*(\vertex-1)}:1.05);
      }
      \draw[thick] (v1)--(v2)--(v3)--(v4)--(v5)--(v6)--cycle;
      \foreach \vertex in {1,...,6} {
        \fill (v\vertex) circle (1.4pt);
        \node[fill=exercisebackground,inner sep=1pt]
          at ({60*(\vertex-1)}:1.33) {$\vertex$};
      }
      \ifnum\panel=1
        \node[courseblue] at (0,-1.75) {$x$: $60^\circ$ rotation};
      \fi
      \ifnum\panel=2
        \node[courseblue] at (0,-1.75) {$x^2$: $120^\circ$ rotation};
      \fi
      \ifnum\panel=3
        \node[ForestGreen] at (0,-1.75) {$y$: reflection\\vertex axis};
      \fi
      \ifnum\panel=4
        \node[BrickRed] at (0,-1.75) {$xy$: reflection\\edge-midpoint axis};
      \fi
    \end{scope}
  }
\end{tikzpicture}
\end{center}

In the homework, you will work out this example in more detail. 


\end{example}

\section{Back to correspondence theorem example}

Recall our  surjective homomorphism
$\varphi:S_4\to S_3$ induced by the action of $S_4$ on the three pairings
\[
  \Pi_1=\{\{1,2\},\{3,4\}\},\qquad
  \Pi_2=\{\{1,3\},\{2,4\}\},\qquad
  \Pi_3=\{\{1,4\},\{2,3\}\}.
\]
Its kernel is the Klein four group
\[
  K=\{1,(1\,2)(3\,4),(1\,3)(2\,4),(1\,4)(2\,3)\}.
\]
We also retain the square symmetries $x=(1\,2\,3\,4)$ and $r=(2\,4)$.

Recall that we have written  $S_3=\{1,a,a^2,b,ab,a^2b\}$, where $a^3=b^2=1$ and
$bab^{-1}=a^{-1}$. Its subgroups are
\[
  \{1\},\quad \langle a\rangle,\quad
  \langle b\rangle,\quad \langle ab\rangle,\quad
  \langle a^2b\rangle,\quad S_3.
\]
Concretely, take $a=(1\,2\,3)$ and $b=(1\,3)$. Then
$a^2=(1\,3\,2)$, $ab=(2\,3)$, and $a^2b=(1\,2)$, with permutation
products composed from right to left. The subgroup inclusion graph is
shown below; inclusions run upward.
\begin{center}
\begin{tikzpicture}[x=1cm,y=1cm,every node/.style={inner sep=4pt,align=center}]
  \node (top) at (0,3.2) {$S_3$};
  \node (cyclic) at (-4.5,1.6) {$\langle a\rangle$\\$\{1,a,a^2\}$};
  \node (b) at (-1.5,1.6) {$\langle b\rangle$\\$\{1,b\}$};
  \node (ab) at (1.5,1.6) {$\langle ab\rangle$\\$\{1,ab\}$};
  \node (aab) at (4.5,1.6) {$\langle a^2b\rangle$\\$\{1,a^2b\}$};
  \node (bottom) at (0,0) {$\{1\}$};
  \draw[courseblue,thick] (bottom) -- (cyclic) -- (top);
  \draw[courseblue,thick] (bottom) -- (b) -- (top);
  \draw[courseblue,thick] (bottom) -- (ab) -- (top);
  \draw[courseblue,thick] (bottom) -- (aab) -- (top);
\end{tikzpicture}
\end{center}

Consequently, there are exactly six subgroups of $S_4$ containing $K$.
Their orders follow from $|\varphi^{-1}(\mathcal H)|=4|\mathcal H|$:
\[
\begin{array}{c|c|c}
  \mathcal H\leq S_3 & \varphi^{-1}(\mathcal H)\leq S_4
    & |\varphi^{-1}(\mathcal H)|\\ \hline
  \{1\} & K & 4\\
  \langle a\rangle=\langle(1\,2\,3)\rangle & A & 12\\
  \langle b\rangle=\langle(1\,3)\rangle & H_1 & 8\\
  \langle ab\rangle=\langle(2\,3)\rangle & H_2 & 8\\
  \langle a^2b\rangle=\langle(1\,2)\rangle & H_3 & 8\\
  S_3 & S_4 & 24
\end{array}
\]

What are the groups $A,H_1,H_2,H_3$? we already identified $A = A_4$. Let's now work with \[ H_1 = \varphi^{-1}( \{ 1, (13) \}). \] We know that $|H_1|=8$. Any guesses as to what it is?

Note that: $\varphi(x)=\varphi(r)=(1\,3)$, so
$\langle x,r\rangle\leq\varphi^{-1}(\langle(1\,3)\rangle)$.
The preimage has order $|K|\cdot2=8$ by the correspondence theorem,
and hence $\langle x,r\rangle=\varphi^{-1}(\langle(1\,3)\rangle)=H_1$.

For $H_2$ and $H_3$, relabel the vertices as
shown below. The opposite vertex pairs now form $\Pi_1$ and $\Pi_3$,
respectively. The curved arrows indicate counterclockwise quarter-turns;
the green dashed lines are diagonal reflection axes.
\begin{center}
\begin{tikzpicture}[x=1cm,y=1cm,every node/.style={font=\small,align=center}]
  \begin{scope}
    \node[courseblue,font=\bfseries] at (0,2.4) {$H_2$};
    \draw[-stealth,courseblue,line width=1pt]
      (45:1.65) arc[start angle=45,end angle=135,radius=1.65];
    \node[courseblue] at (0,1.9) {$x_2$};
    \draw[dashed,ForestGreen,line width=1pt] (-1.3,-1.3)--(1.3,1.3);
    \node[ForestGreen,right] at (1.35,1.3) {$r_2$};
    \draw[thick] (-0.9,-0.9)--(-0.9,0.9)--(0.9,0.9)--(0.9,-0.9)--cycle;
    \fill (0.9,0.9) circle (1.5pt) node[above] {$1$};
    \fill (-0.9,0.9) circle (1.5pt) node[above] {$3$};
    \fill (-0.9,-0.9) circle (1.5pt) node[below] {$2$};
    \fill (0.9,-0.9) circle (1.5pt) node[below] {$4$};
    \node[courseblue] at (0,-1.45) {opposite pairs: $\Pi_1$};
    \node at (0,-2) {$x_2=(1\,3\,2\,4),\quad r_2=(3\,4)$};
  \end{scope}
  \begin{scope}[xshift=7cm]
    \node[BrickRed,font=\bfseries] at (0,2.4) {$H_3$};
    \draw[-stealth,BrickRed,line width=1pt]
      (45:1.65) arc[start angle=45,end angle=135,radius=1.65];
    \node[BrickRed] at (0,1.9) {$x_3$};
    \draw[dashed,ForestGreen,line width=1pt] (-1.3,-1.3)--(1.3,1.3);
    \node[ForestGreen,right] at (1.35,1.3) {$r_3$};
    \draw[thick] (-0.9,-0.9)--(-0.9,0.9)--(0.9,0.9)--(0.9,-0.9)--cycle;
    \fill (0.9,0.9) circle (1.5pt) node[above] {$1$};
    \fill (-0.9,0.9) circle (1.5pt) node[above] {$2$};
    \fill (-0.9,-0.9) circle (1.5pt) node[below] {$4$};
    \fill (0.9,-0.9) circle (1.5pt) node[below] {$3$};
    \node[BrickRed] at (0,-1.45) {opposite pairs: $\Pi_3$};
    \node at (0,-2) {$x_3=(1\,2\,4\,3),\quad r_3=(2\,3)$};
  \end{scope}
\end{tikzpicture}
\end{center}
The induced permutations of the three pairings are
\[
  \varphi(x_2)=\varphi(r_2)=(2\,3)=ab,\qquad
  \varphi(x_3)=\varphi(r_3)=(1\,2)=a^2b.
\]
Each square symmetry group has order $8$, as does the corresponding
preimage. Therefore
\[
\begin{aligned}
  H_2&=\langle x_2,r_2\rangle=\varphi^{-1}(\langle ab\rangle),\\
  H_3&=\langle x_3,r_3\rangle=\varphi^{-1}(\langle a^2b\rangle).
\end{aligned}
\]
Both are copies of $D_8$ inside $S_4$ containing $K$.

\begin{exercise}
Use the correspondence theorem to explain why none of the three
subgroups $H_1,H_2,H_3$ is normal in $S_4$.
\end{exercise}

\section{External direct products}

There are two ways to approach products of groups:
\begin{itemize}
  \item An \emph{external} product starts with groups $G_1,G_2$ and
  constructs a new group $G_1\times G_2$.
  \item An \emph{internal} product starts with a group $G$ and asks
  whether suitable subgroups $G_1,G_2\leq G$ give an isomorphism
  $G\cong G_1\times G_2$.
\end{itemize}
We first develop the external construction, then explain how to recognize
a direct product inside a given group.

\begin{define}[External direct product]
Let $G_1$ and $G_2$ be groups. Their \emph{direct product}
$G_1\times G_2$ is the Cartesian product of their underlying sets,
with multiplication defined componentwise:
\[
  (g_1,g_2)(g'_1,g'_2)=(g_1g'_1,g_2g'_2).
\]
\end{define}

\begin{lemma}\label{lem:product-group}
The direct product is a group. Its identity is $(e_1,e_2)$ and
\[
  (g_1,g_2)^{-1}=(g_1^{-1},g_2^{-1}).
\]
\end{lemma}

\begin{exercise}
Verify the group axioms for $G_1\times G_2$, using the corresponding
axioms in each factor. If both groups are finite, show that
$|G_1\times G_2|=|G_1|\,|G_2|$.
\end{exercise}

\subsection{The factors as normal subgroups}

The maps
\[
  i_1:G_1\longrightarrow G_1\times G_2,\quad g_1\longmapsto(g_1,e_2),
\]
\[
  i_2:G_2\longrightarrow G_1\times G_2,\quad g_2\longmapsto(e_1,g_2)
\]
are injective homomorphisms. We often identify the factors with their
images, but the subgroups inside the product are, more precisely,
\[
  \widetilde G_1=G_1\times\{e_2\},\qquad
  \widetilde G_2=\{e_1\}\times G_2.
\]

\begin{lemma}
Both $\widetilde G_1$ and $\widetilde G_2$ are normal in
$G_1\times G_2$, and
\[
  \widetilde G_1\cap\widetilde G_2=\{(e_1,e_2)\}.
\]
\end{lemma}
\begin{proof}
For $(g_1,g_2)\in G_1\times G_2$ and $h_1\in G_1$,
\[
  (g_1,g_2)(h_1,e_2)(g_1,g_2)^{-1}
  =(g_1h_1g_1^{-1},e_2)\in\widetilde G_1.
\]
Thus $\widetilde G_1$ is normal; the same computation in the second
coordinate proves normality of $\widetilde G_2$.
An element in both subgroups must have first coordinate $e_1$ and
second coordinate $e_2$, proving the intersection statement.
\end{proof}

\section{Products of cyclic groups}

Write $C_d$ for a cyclic group of order $d$.

\begin{lemma}\label{lem:coprime-cyclic}
If $r,s\geq1$ and $\gcd(r,s)=1$, then
\[
  C_{rs}\cong C_r\times C_s.
\]
\end{lemma}

\begin{warningbox}
The coprimality assumption matters. For example,
$C_4\not\cong C_2\times C_2$. Every element $(u,v)$ of
$C_2\times C_2$ satisfies
\[
  (u,v)^2=(u^2,v^2)=(1,1),
\]
whereas $C_4$ has an element of order $4$. An isomorphism preserves
element orders, so these groups cannot be isomorphic.
\end{warningbox}

\begin{proof}[Proof of Lemma~\ref{lem:coprime-cyclic}]
Choose generators
\[
  C_{rs}=\langle x\rangle,\qquad C_r=\langle z\rangle,
  \qquad C_s=\langle y\rangle.
\]
Define
\[
  \varphi:C_{rs}\to C_r,\quad x^i\mapsto z^i,
  \qquad
  \psi:C_{rs}\to C_s,\quad x^i\mapsto y^i.
\]
These maps are well defined: if $x^i=x^j$, then $rs\mid(i-j)$,
so both $r$ and $s$ divide $i-j$, giving $z^i=z^j$ and $y^i=y^j$.
They are homomorphisms, since, for instance,
\[
  \varphi(x^ix^j)=z^{i+j}=z^iz^j.
\]
Define a map using these two coordinates:
\[
  \Phi:C_{rs}\longrightarrow C_r\times C_s,
  \qquad x^i\longmapsto(z^i,y^i).
\]
It is a homomorphism because multiplication is componentwise:
\[
  \Phi(x^ix^j)=(z^{i+j},y^{i+j})
  =(z^i,y^i)(z^j,y^j)=\Phi(x^i)\Phi(x^j).
\]
The two coordinate kernels are
\[
  \ker(\varphi)=\langle x^r\rangle,\qquad
  \ker(\psi)=\langle x^s\rangle.
\]
Consequently, $x^i\in\ker(\Phi)$ if and only if $r\mid i$ and
$s\mid i$. Since $\gcd(r,s)=1$, this is equivalent to $rs\mid i$,
so $\ker(\Phi)=\{1\}$ and $\Phi$ is injective.

Both its domain and its codomain have order $rs$. Thus $\Phi$ is also
surjective, and hence an isomorphism. Equivalently, the counting
formula gives $|\im(\Phi)|=rs$, so its image is all of
$C_r\times C_s$.
\end{proof}

\section{Internal direct products}

\subsection{The product of two subgroups}

Let $H,K\leq G$. The \emph{product set} of $H$ and $K$ is
\[
  HK=\{hk\mid h\in H,\ k\in K\}\subseteq G.
\]
Here $HK$ consists of elements of $G$, whereas $H\times K$ consists
of ordered pairs and has the external direct product group structure.

\begin{warningbox}
The set $HK$ need not be a subgroup of $G$.
\end{warningbox}

\begin{example}
In $S_3$, let $H=\langle(1\,2)\rangle$ and
$K=\langle(2\,3)\rangle$. Composing permutations from right to left,
\[
  HK=\{1,(1\,2),(2\,3),(1\,2\,3)\}.
\]
This set is not closed under multiplication: both $(2\,3)$ and
$(1\,2)$ belong to $HK$, but their product $(1\,3\,2)$ does not.
\end{example}

\begin{lemma}[Unique factorization]\label{lem:unique-factorization}
If $H\cap K=\{1\}$, the multiplication map of sets
\[
  p:H\times K\longrightarrow G,\qquad (h,k)\longmapsto hk
\]
is injective. Thus each element of $HK$ has a unique expression $hk$
with $h\in H$ and $k\in K$.
\end{lemma}
\begin{proof}
If $h_1k_1=h_2k_2$, then
\[
  h_2^{-1}h_1=k_2k_1^{-1}.
\]
The left side lies in $H$ and the right side lies in $K$. Their common
value therefore belongs to $H\cap K=\{1\}$, so $h_1=h_2$ and $k_1=k_2$.
\end{proof}

\begin{lemma}\label{lem:product-subgroup}
If either $H$ or $K$ is normal in $G$, then $HK=KH$ and $HK\leq G$.
\end{lemma}
\begin{proof}
Suppose first that $H\trianglelefteq G$. For $h\in H$ and $k\in K$,
\[
  hk=k(k^{-1}hk)\in KH,\qquad
  kh=(khk^{-1})k\in HK,
\]
because the conjugates $k^{-1}hk$ and $khk^{-1}$ lie in $H$.
Thus $HK=KH$.

The identity lies in $HK$. For closure under multiplication, write
$k_1h_2=h'_2k'_1$ using $KH=HK$. Then
\[
  (h_1k_1)(h_2k_2)=(h_1h'_2)(k'_1k_2)\in HK.
\]
Also, $(hk)^{-1}=k^{-1}h^{-1}\in KH=HK$, so $HK$ is closed under
inverses and hence is a subgroup. If $K\trianglelefteq G$, interchange
$H$ and $K$ in the argument.
\end{proof}

\Needspace{10\baselineskip}
\subsection{Recognizing a direct product}

\begin{define}[Internal direct product]
A group $G$ is the \emph{internal direct product} of subgroups $H$ and
$K$ if
\[
  H\trianglelefteq G,\qquad K\trianglelefteq G,\qquad
  H\cap K=\{1\},\qquad HK=G.
\]
\end{define}

\begin{theorem}[Recognizing direct products]\label{thm:internal-product}
If $G$ is the internal direct product of $H$ and $K$, then
\[
  G\cong H\times K.
\]
The isomorphism sends $g=hk$ to $(h,k)$.
\end{theorem}
\begin{proof}
Since $HK=G$, every $g\in G$ can be written $g=hk$. By
Lemma~\ref{lem:unique-factorization}, this expression is unique. Hence
\[
  \pi:G\longrightarrow H\times K,\qquad hk\longmapsto(h,k)
\]
is well defined.

\emph{First we show that elements of $H$ commute with elements of $K$.}
For $h\in H$ and $k\in K$, consider the commutator
\[
  [h,k]=hkh^{-1}k^{-1}.
\]
Normality of $K$ gives $hkh^{-1}\in K$, so $[h,k]\in K$.
Normality of $H$ gives $kh^{-1}k^{-1}\in H$, so
$[h,k]=h(kh^{-1}k^{-1})\in H$. Consequently,
\[
  [h,k]\in H\cap K=\{1\},\qquad\text{and thus }hk=kh.
\]

Now let $g_1=h_1k_1$ and $g_2=h_2k_2$. Since $k_1h_2=h_2k_1$,
\[
  g_1g_2=h_1k_1h_2k_2=(h_1h_2)(k_1k_2).
\]
It follows that
\[
  \pi(g_1g_2)=(h_1h_2,k_1k_2)
  =(h_1,k_1)(h_2,k_2)=\pi(g_1)\pi(g_2).
\]
Thus $\pi$ is a homomorphism. It is bijective, with inverse
$p(h,k)=hk$, and hence is an isomorphism.
\end{proof}

\begin{exercise}
Verify directly that the maps $\pi$ and $p$ in
Theorem~\ref{thm:internal-product} are inverses. Then check that the
subgroups $\widetilde G_1$ and $\widetilde G_2$ of the external product
$G_1\times G_2$ satisfy all four conditions in the definition of an
internal direct product.
\end{exercise}

\begin{remark}
Both normality assumptions matter. If only $H\trianglelefteq G$ is
assumed, together with $H\cap K=\{1\}$ and $HK=G$, the resulting
structure is called a \emph{semidirect product}, written
$G\cong H\rtimes K$; the factors need not commute.

For example, in $S_3$ take $H=\langle a\rangle$ with $a=(1\,2\,3)$
and $K=\langle b\rangle$ with $b=(1\,2)$. Then $H\trianglelefteq S_3$,
$H\cap K=\{1\}$, and
\[
  HK=\{1,a,a^2,b,ab,a^2b\}=S_3.
\]
But $K$ is not normal: $aba^{-1}=(2\,3)\notin K$.
Moreover, $bab^{-1}=a^{-1}\ne a$, so the factors do not commute.
Thus $S_3\cong C_3\rtimes C_2$ is a semidirect product; it is not
isomorphic to the abelian group $C_3\times C_2$.
\end{remark}

\end{document}
